\section{Cycle Finding}
\label{sec:cses-cycle-finding}

\problemstatement{Given a directed graph of $n$ nodes and $m$ edges with
(possibly negative) weights, find a negative cycle and print the nodes on
it, or output \texttt{NO CYCLE}.}

\begin{patternbox}
Detecting a negative cycle is the failure mode of \textbf{Bellman--Ford}:
if any edge can still be relaxed after $n-1$ rounds, a negative cycle is
reachable. Following predecessor pointers back $n$ steps lands on the cycle,
which is then unwound.
\end{patternbox}

\subsection*{Relax $n$ times, then trace back}

Initialise all distances to $0$ (so every node is a potential start) and
relax all edges $n$ times, recording a predecessor on each successful
relaxation. If some edge relaxes on the $n$-th round, its endpoint $x$ is
affected by a negative cycle. Stepping from $x$ back through predecessors
$n$ times guarantees landing on a node $y$ that is truly on the cycle;
following predecessors from $y$ until $y$ repeats reconstructs the cycle.

\begin{ideabox}[The n-th Relaxation Exposes a Negative Cycle]
Shortest paths use at most $n-1$ edges, so a further improvement can only
come from a cycle with negative total weight. Walking back $n$ predecessors
from the improved node skips the ``tail'' into the cycle, ensuring the
printed sequence is the cycle itself.
\end{ideabox}

\subsection*{Algorithm}

\begin{enumerate}
  \item $\textit{dist}[i]=0$ for all $i$; relax all edges $n$ times,
        storing predecessors; remember the last node improved on round
        $n$.
  \item If none improved, print \texttt{NO CYCLE}.
  \item Else step back $n$ times to reach a cycle node, then collect the
        cycle by following predecessors until it repeats; print it.
\end{enumerate}

\subsection*{Dry run}

\dryrun{$1\!\to\!2(1)$, $2\!\to\!3(-3)$, $3\!\to\!1(1)$ (cycle sum
$-1$).}

\begin{itemize}
  \item After repeated relaxation, an edge keeps improving -- a negative
        cycle exists.
  \item Tracing back yields the cycle $1\to2\to3\to1$.
\end{itemize}

\subsection*{C++ implementation}

\begin{lstlisting}
#include <bits/stdc++.h>
using namespace std;

int main() {
    int n, m;
    cin >> n >> m;
    vector<array<long long,3>> edges(m);
    for (auto& e : edges) cin >> e[0] >> e[1] >> e[2];

    vector<long long> dist(n + 1, 0);
    vector<int> pre(n + 1, -1);
    int x = -1;
    for (int i = 0; i < n; i++) {
        x = -1;
        for (auto& [u, v, w] : edges)
            if (dist[u] + w < dist[v]) {
                dist[v] = dist[u] + w;
                pre[v] = u;
                x = v;                            // node improved this round
            }
    }

    if (x == -1) { cout << "NO CYCLE\n"; return 0; }
    for (int i = 0; i < n; i++) x = pre[x];       // step into the cycle

    vector<int> cyc;
    for (int v = x; ; v = pre[v]) {
        cyc.push_back(v);
        if (v == x && cyc.size() > 1) break;
    }
    reverse(cyc.begin(), cyc.end());
    cout << "YES\n";
    for (int v : cyc) cout << v << " ";
    cout << "\n";
    return 0;
}
\end{lstlisting}

\begin{complexitybox}
\textbf{Time Complexity.} $O(nm)$ -- Bellman--Ford relaxation.
\textbf{Space Complexity.} $O(n+m)$.
\end{complexitybox}

\begin{takeawaybox}
Negative-cycle finding is Bellman--Ford relaxed one extra round, then a
back-walk of $n$ predecessors to land firmly on the cycle. Initialising all
distances to $0$ makes every node a virtual source, so any negative cycle
anywhere is caught.
\end{takeawaybox}
